\documentclass[12pt]{amsart}
\usepackage{graphicx} % Required for inserting images
\usepackage{amsmath,amsthm,amsfonts,bbm,bm,color,cite,enumerate}
\usepackage[foot]{amsaddr}
\newtheorem{theorem}{Theorem}[section]
\newtheorem{cor}[theorem]{Corollary}
\newtheorem{lem}[theorem]{Lemma}
\newtheorem{pro}[theorem]{Proposition}
\newtheorem{conj}[theorem]{Conjecture}
%\newtheorem{assumption}[theorem]{Assumption}
\newtheorem{ex}[theorem]{Example}
\theoremstyle{definition}
\newtheorem{de}[theorem]{Definition}
\theoremstyle{remark}
\newtheorem{rem}[theorem]{Remark}
\numberwithin{equation}{section}
\usepackage[margin=1in]{geometry}
\usepackage[nameinlink,noabbrev]{cleveref}

\newtheorem{assumption}{Assumption}
\crefname{assumption}{assumption}{assumptions}

\title[NLS numerical]{Numerical study of a energy cascade model derived from dispersive equation}
\author{Jayleen Jiang}
\author{Yao Li}
\thanks{Department of Mathematics and Statistics, University of Massachusetts Amherst, Amherst, Massachusetts, the United States, 01003}
\email{liyao@umass.edu}

\date{\today}

\begin{document}
\maketitle

\section{Introduction}
%skip it for now
% derivation of energy cascade model, state the equation here


% significance and challenge
% known result

\section{Preliminary}
% (long chain) nonequilibrium thermodynamics: Thermal conductivity, Fourier's law, entropy production rate, Fluctuation theorem

% Fokker-Planck equation, eigenvalue and eigenfunction decomposition

% How to solve Fokker-Planck, how to find eigenfunction

\section{Statistical physics properties of long cascade chain}
\subsection{Nonequilibrium steady state}
% 1.state how we did the simulation. How we confirm that the simulation is stabilized (expectation does not change much with the increasing sample size). What is the energy profile versus N for the "closed chain" (case 0)? What is the energy profile versus N for the "open chain" (case 4)

% 2.terminal energy of open chain versus N. When $N$ gets larger in the "open chain", what is the "terminal energy"? Plot the "terminal energy" versus increasing N. (TODO:)Add some interpretation (connection between $\Lambda_n$ and the actual Fourier modes). Plot the location of the first N groups on a 2D plane. 


% 3.local thermodynamic equilibrium. In a long chain (N = 25?), find the marginal probability density function of (I_12, I_13, theta_12), denoted by $\rho_{middle}$. Is $\log (\rho_{middle}) $  proportional to 
%$$
%-(I_{12}^2 + I_{13}^2)/4 + I_{12}I_{13} \cos \theta_{12}
%$$
%??


1, Compare marginal PDF of $(I_a, I_{a+1}, \theta_a)$ and the 2-mode Gibbs by comparing $-\log(PDF)$ and $\frac{1}{2} (I_1 + I_2)^2 - \frac{1}{4} (I_1^2 + I_2^2) + I_1I_2\cos \theta_1$. (They are different)


2, Compare marginal PDF of $(I_a, I_{a+1}, \theta_a)$ of NESS and the same marginal PDF of equilibrium by comparing the $-\log (PDF)$ of both. 

Equilibrium: $p_\beta(I_a, I_{a+1}, \theta_a) : = \int_A \exp(-H/\beta) \mathrm{d}x$, where $A$ is the collection of variables other than $(I_a, I_{a+1}, \theta_a)$.

Nonequilibrium: Does the marginal PDF of $(I_a, I_{a+1}, \theta_a)$, denoted by $q(I_a, I_{a+1}, \theta_a)$ approach to $p_{\beta_0}(I_a, I_{a+1}, \theta_a)$ of some $\beta_0$ as the length of the chain ($n$) increases?? How does $\beta_0$ depend on $(T_L, T_R, a/n)$? 



Test: fix $a/n$. Let $n$ increase. Let $q_\beta$ be the equilibrium with boundary temperature $\beta$. Find the minimum of $\min_x\|\exp(x\log q) - p\|$. So the effective boundary temperature is $\beta_0:= \beta/x$. Compare $q_{\beta_0}$ and $p$.



\subsection{Thermal conductivity}
% Energy flux versus length of the chain, versus energy profile gradient 

% Histogram of the energy flux. Is $P[ energy flux > C] \approx exp(-r C)$ ?

\subsection{Entropy production and its fluctuation}

% compute the entropy production rate as in the second law of thermodynamics

% Does it satisfy fluctuation-dissipation theorem??


\section{Detailed property of short chain}
\subsection{Numerical solution of the noneqilibrium steady state}
% Solution of Fokker-Planck

\subsection{Mechansim of stabilization}
% Mechanism of stabilization : high energy and low energy

\subsection{Eigenfunction of the generator}
% Eigenfunction of generator and interpretation ?

\end{document}